% [NEW]: Ket luan - viet moi hoan toan
\chapter*{Kết luận}
\addcontentsline{toc}{chapter}{Kết luận}

Khóa luận này đã nghiên cứu bài toán sinh câu hỏi trắc nghiệm tự động từ ảnh chụp tài liệu giáo dục, sử dụng các mô hình ngôn ngữ lớn cỡ nhỏ (sLLM) có khả năng triển khai tại chỗ với chi phí thấp. Đây là bài toán có ý nghĩa thực tiễn cao trong bối cảnh nhu cầu số hóa giáo dục ngày càng tăng, đặc biệt tại các môi trường hạn chế về tài nguyên tính toán và ngân sách.

% [MODIFIED]: Chinh lai phan tong ket dong gop de phu hop voi kien truc thuc te da mo ta o Chuong 3, bo mo rong DCP cu va them vai tro cua tac tu Hieu chinh.
\textbf{Đóng góp chính.} Khóa luận đề xuất bốn đóng góp cốt lõi. Thứ nhất, kiến trúc đa tác tử, trong đó ba tác tử sinh câu hỏi hoạt động song song theo từng cấp độ Bloom và được bao bọc bởi Tác tử Phân loại, Tác tử Học sinh cùng Tác tử Hiệu chỉnh để tạo thành hai vòng kiểm soát chất lượng độc lập. Tác tử Phân loại áp dụng kỹ thuật chống thiên lệch bằng ánh xạ ID tạm thời để đánh giá cấp độ Bloom một cách khách quan, còn Tác tử Học sinh và Tác tử Hiệu chỉnh phối hợp nhằm loại bỏ các câu hỏi khó giải hoặc có phương án nhiễu kém chất lượng. Thứ hai, khung đánh giá bốn tiêu chí (Solvability, Distractor Quality, Alignment, Clarity) kết hợp phương pháp LLM-as-a-Judge cho phép đánh giá tự động chất lượng câu hỏi trắc nghiệm một cách nhất quán và có hệ thống. Thứ ba, quy trình sinh dữ liệu OCR tổng hợp kết hợp Browser Engineering và Augraphy tạo ra tập dữ liệu ảnh tài liệu đa dạng phục vụ huấn luyện và đánh giá hệ thống. Thứ tư, ứng dụng web Openlet tích hợp đầy đủ kiến trúc nghiên cứu thành sản phẩm demo hoàn chỉnh, cho phép người dùng trải nghiệm toàn bộ quy trình từ chụp ảnh tài liệu đến nhận câu hỏi trắc nghiệm.

\textbf{Kết quả chính.} Thực nghiệm cho thấy kiến trúc đa tác tử cải thiện chất lượng sinh câu hỏi của sLLM từ 30\% đến 90\% trên tất cả các chỉ số đánh giá so với phương pháp sinh trực tiếp, với chi phí suy luận dưới \$0.50 cho mỗi 100 đoạn văn đầu vào. Đáng chú ý, sLLM tốt nhất (Gemma 3 kết hợp kiến trúc đa tác tử) đạt hiệu suất cạnh tranh với các mô hình ngôn ngữ lớn thương mại ở Cấp độ 1 và Cấp độ 2 theo thang Bloom. Tác tử Phân loại đạt điểm Alignment gần 1.0 nhờ kỹ thuật ánh xạ ID tạm thời, khẳng định hiệu quả của phương pháp chống thiên lệch vị trí.

\textbf{Ý nghĩa thực tiễn.} Kết quả nghiên cứu mở ra khả năng ứng dụng thực tế cho lĩnh vực công nghệ giáo dục (EdTech) với chi phí thấp. Giáo viên có thể sử dụng hệ thống để nhanh chóng tạo bộ câu hỏi trắc nghiệm đa dạng về cấp độ nhận thức từ tài liệu giảng dạy, trong khi học sinh có thể tự ôn luyện bằng cách chụp ảnh sách giáo khoa và nhận câu hỏi kiểm tra ngay lập tức.

\textbf{Hướng phát triển.} Trong tương lai, nghiên cứu có thể được mở rộng theo các hướng sau: (1) hỗ trợ tiếng Việt và các ngôn ngữ khác ngoài tiếng Anh; (2) tiến hành đánh giá với giáo viên và chuyên gia giáo dục thực tế để kiểm chứng chất lượng câu hỏi; (3) fine-tuning sLLM trên tập dữ liệu câu hỏi trắc nghiệm chuyên biệt để cải thiện chất lượng ở các cấp độ Bloom cao; và (4) tích hợp hệ thống với các nền tảng quản lý học tập (LMS) phổ biến để đưa giải pháp đến gần hơn với người dùng cuối.
